\section{Discussion and open questions}

The consistency criterion in \cref{obj:thm:consistency-threshold} links sample size, categorical support, and overlap through a single joint requirement. With \(n\) observed units, \(d\) covariate cells, overlap floor \(\epsilon\), and logarithmic alphabet scale \(L_d\), consistency requires
\[
\frac{d}{n\epsilon L_d}\longrightarrow 0.
\]
Thus, sample growth must compensate jointly for an expanding covariate alphabet and declining overlap. The criterion makes this tradeoff explicit: the relevant comparison is between the number of cells and sample size multiplied by the overlap floor and logarithmic alphabet scale.

For categorical treatment-choice problems, the minimax comparison in \cref{obj:thm:matched-frontier} provides a benchmark for assessing the precision of optimal-value estimation. In the nonsaturated regime, the squared-error scale is \(d/(n\epsilon L_d)\), up to universal comparison constants. This scale supports assessments of how changes in sample size, covariate resolution, and overlap affect the smallest achievable worst-case risk. The guarantee in \cref{obj:thm:observable-upper} supplies an observed-data estimator attaining that scale, linking the precision benchmark to the construction in \cref{obj:def:armwise-estimator}.

The benchmark also accommodates heterogeneity in treatment assignment and indifference between treatment arms. Within the class in \cref{obj:def:observed-class}, propensities may differ across occupied cells while satisfying the public overlap bounds. The value in \cref{obj:def:observed-value} uses the larger conditional outcome mean in each cell, including cells where the two means coincide. Such ties leave the value well defined even when either treatment attains the same cellwise optimum. Under the consistency and conditional exchangeability conditions in \cref{obj:ass:consistency,obj:ass:conditional-exchangeability}, \cref{obj:prop:identification} gives this observed-law benchmark its causal interpretation.

\subsection{Limitations and future work}

The attaining estimator is explicit as a measurable finite-sample statistic, but its displayed conditional average enumerates \(2^k\) mark assignments for each auxiliary count \(k\le n\), in addition to evaluating four-dimensional Jackson integrals or precomputed polynomial coefficients. The result therefore establishes statistical attainment rather than an efficient numerical algorithm. Finding a polynomial-time approximation with the same uniform risk guarantee, and quantifying the numerical error permitted without changing the rate, remain open computational questions.

The proved rate comparison determines the order of minimax risk. Leading-constant convergence is a separate question, particularly when the covariate alphabet grows and overlap vanishes together. The following question isolates that issue within the nonsaturated regime.

\begin{remarkv}{oeq:sharp-path-constant}[Open questions on path constants]
Consider admissible nonsaturated sequences satisfying
\begin{itemize}
\item \textbf{(Growing alphabet.)} \(d\to\infty\).
\item \textbf{(Vanishing overlap.)} \(\epsilon\to0\).
\item \textbf{(Nonsaturation.)} \(d/(n\epsilon L_d)\to0\), where \(L_d\) is the logarithmic alphabet scale.
\end{itemize}
For \(\mathfrak R_{n,d,\epsilon}\), the minimax squared-error risk, consider the normalized sequence
\[
\frac{n\epsilon L_d}{d}\,\mathfrak R_{n,d,\epsilon}.
\]
It remains open whether there exists a single finite positive constant to which the entire normalized-risk sequence converges along every admissible path. Alternatively, do there exist two admissible overlap paths along which the entire normalized-risk sequences converge to distinct finite positive constants? Failure of convergence along an admissible path is a further possibility.

Future work could investigate a cone-constrained approximation--moment saddle for the boundary-propensity likelihood and an observable no-loss transfer.
\end{remarkv}

The nonsaturation condition places these sequences in the region where the minimax risk tends to zero. Along them, \cref{obj:thm:matched-frontier} bounds the normalized risk above and below by universal positive constants. Convergence to a common constant would sharpen that comparison into a path-independent calibration of leading risk. Distinct limits along two admissible paths would instead reveal dependence on how sample size, support, and overlap evolve jointly. The alternatives in \cref{obj:oeq:sharp-path-constant} concern convergence of entire sequences; convergence along selected subsequences would leave the stated question unresolved.

The proposed approximation--moment direction seeks a sharp characterization compatible with the boundary-propensity likelihood and the constrained prior pairs in \cref{obj:def:constrained-prior-separation}. An observable no-loss transfer would aim to carry the resulting leading constant into the observable statistical problem without deterioration. Developing these two steps could connect sharp approximation and moment calculations to a leading-risk calibration beyond the established rate comparison.
